\chapter{Design of the Proposed Model}
\label{chapter:design}

This chapter describes the design of the proposed model, referred to as \modelname in what follows, at the
level of its components, the data that passes between them and the algorithms they execute. Its
realisation in a particular library is left to Chapter~\ref{chapter:implementation}. \modelname extends
ViP3D \cite{Gu2023ViP3D} and reuses most of its pipeline, so the chapter first describes the parts of ViP3D
that it reuses and then the parts that this thesis replaces or adds, which follow the starting points of
Section~\ref{sec:summary}. It closes with the training schedule and an estimate of the time and memory
complexity of every stage.


% ============================================================================
\section{Overview}
\label{sec:design-overview}
% ============================================================================

Our model processes a scene frame by frame and carries a set of agent queries from one frame to the next.
One frame passes through these 5 stages.

\begin{enumerate}
    \item \textbf{Feature extraction.} The point cloud is encoded into a \ac{bev} feature map and, when the
    camera is used, the six images into multi-level image features.
    \item \textbf{Query initialisation.} A class-wise heatmap is predicted over the \ac{bev} map, and its
    strongest peaks initialise the query slots that do not follow an object.
    \item \textbf{Decoding.} A two-layer decoder refines all queries, the first layer by reading the
    \ac{bev} map and the second by attending to the camera features, and after each layer predicts a box,
    a class and a velocity for every query.
    \item \textbf{Tracking update.} Queries are assigned identities, their embeddings are refined by a
    memory of their own past and by interaction with the other tracks, and the reference points of the
    surviving tracks are moved into the coordinate frame of the next frame.
    \item \textbf{Trajectory prediction.} The embeddings of the tracked agents, together with the lanes of
    an \ac{hdmap}, are decoded into several future trajectories per agent.
\end{enumerate}

The camera part of stage 1 and stages 4 and 5 are taken from ViP3D. The \ac{lidar} part of stage 1 is
added, stage 2 replaces the learned query embedding of ViP3D, and stage 3 replaces its six-layer camera
decoder.

%figure: Block scheme of the complete \modelname architecture. Two input branches at the top: the point cloud
% passes the LiDAR branch and becomes the BEV feature map, the six camera images pass ResNet-50 and the FPN and
% become multi-level image features. The BEV map feeds the class-wise heatmap head, whose strongest peaks become
% the queries of the empty slots, joined by the track queries carried from the previous frame. All queries pass
% decoder layer 0 (BEV read) and decoder layer 1 (SMCA camera attention), each followed by its prediction branches.
% The output goes to the memory bank and the query interaction module, which return the track queries for the
% next frame, and the embeddings of tracked agents go with the HD-map lanes to the trajectory predictor.
% Mark the parts added by this thesis (heatmap head, query initialisation, two-layer decoder) in a distinct colour.
% The figure presentation/figures/initvip_architecture.png can serve as a starting point.

The symbols used throughout the chapter are listed in Table~\ref{tab:design-notation}, together with the
values they take in the configuration used for the experiments.

% TODO - maybe remove the concrete values and move them to the implementation chapter
\begin{table}[htbp]
    \centering
    \caption{Symbols used in this chapter and their values in the configuration used.}
    \label{tab:design-notation}
    \begin{tabular}{llr}
        \toprule
        Symbol & Meaning & Value \\
        \midrule
        $N_q$ & empty query slots appended in every frame & 300 \\
        $N_a$ & track queries carried over from the previous frame & varies \\
        $Q = N_q + N_a$ & queries processed by the decoder in one frame & varies \\
        $D$ & embedding width & 256 \\
        $N_{\mathrm{cls}}$ & object classes & 7 \\
        $H \times W$ & \ac{bev} grid & $256 \times 256$ \\
        $s$ & \ac{bev} cell size & 0.4\,m \\
        $C$ & channels of the \ac{bev} feature map & 384 \\
        $N_{\mathrm{cam}}$ & cameras & 6 \\
        $h \times w$ & camera feature map read by the fusion layer & $68 \times 120$ \\
        $T$ & frames in a training clip & 3 \\
        $L$ & length of the memory bank & 4 \\
        $M$ & predicted trajectories per agent & 6 \\
        $T_f$ & predicted future positions, 6\,s at 2\,Hz & 12 \\
        \bottomrule
    \end{tabular}
\end{table}


% ============================================================================
\section{Components Reused from ViP3D}
\label{sec:design-vip3d}
% ============================================================================

This section describes the parts of ViP3D that \modelname reuses, namely the agent queries, the camera
feature extraction, the prediction branches, the tracking modules and the trajectory predictor. The general
idea of ViP3D was introduced in Section~\ref{sec:rw-decoding}, so the focus here is on how these parts
operate and exchange data.

\subsection{Agent Queries and Track State}
\label{subsec:design-state}

The central object of ViP3D is the agent query. Following MOTR \cite{Zeng2022MOTR}, a frame is processed
with two kinds of queries. Track queries are carried over from the previous frame and follow objects that
have already been found, while a fixed number $N_q$ of empty slots is appended in every frame to find
objects that are new. The decoder therefore processes $Q = N_q + N_a$ queries in a frame, where $N_a$ is the
number of tracks alive at that moment. Each slot $i$ holds a state

\begin{equation}
    \mathcal{S}_i = \left( \mathbf{p}_i,\, \mathbf{c}_i,\, \mathbf{r}_i,\, \mathbf{z}_i,\, \mathbf{h}_i,\,
    \mathrm{id}_i,\, s_i,\, \mathbf{B}_i \right),
    \label{eq:design-state}
\end{equation}

where $\mathbf{p}_i \in \mathbb{R}^D$ and $\mathbf{c}_i \in \mathbb{R}^D$ are the positional and the content
half of the query, $\mathbf{r}_i \in \mathbb{R}^3$ is its reference point, $\mathbf{z}_i \in \mathbb{R}^3$ is
the logarithm of the width, length and height of its box, $\mathbf{h}_i \in \mathbb{R}^D$ is the output
embedding of the last decoder layer, $\mathrm{id}_i$ is the identity of the object the slot follows or $-1$
for an empty slot, $s_i$ is its score, and $\mathbf{B}_i \in \mathbb{R}^{L \times D}$ is a memory bank entry of its
past embeddings.

The reference point is not stored in metres. Its coordinates are first normalised to the point cloud range
$[x_{\min}, x_{\max}] \times [y_{\min}, y_{\max}] \times [z_{\min}, z_{\max}]$ and then passed through the
inverse of the sigmoid function $\sigma$,

\begin{equation}
    \mathbf{r}_i = \sigma^{-1}\!\left( \tilde{\mathbf{r}}_i \right), \qquad
    \tilde{r}_{i,x} = \frac{x_i - x_{\min}}{x_{\max} - x_{\min}},
    \label{eq:design-invsig}
\end{equation}

and likewise for $y$ and $z$. Every refinement of the reference point is then an unbounded offset added to
$\mathbf{r}_i$, and applying $\sigma$ afterwards maps the result back into the range, so that a refined point
can never leave it.

%figure: Sequence diagram of the query slots over three consecutive frames. In every frame, a column of track
% queries carried from the previous frame (coloured by identity) stands next to a column of N_q empty slots.
% After decoding, some empty slots become new tracks (birth), some tracks end (death), and the survivors move on
% to the next frame, where N_q fresh empty slots are appended again. The figure should make visible that the
% number of queries in a frame is N_q + N_a.

\subsection{Camera Feature Extraction}
\label{subsec:design-camera}

Each of the $N_{\mathrm{cam}}$ images is encoded by a ResNet-50 \cite{He2016ResNet}, whose last two stages
use deformable convolutions \cite{Zhu2019DCNv2}, followed by an \ac{fpn} \cite{Lin2017FPN} that produces four
feature maps of $D$ channels at strides 8, 16, 32 and 64 of the input image. Before the decoder reads them,
each map receives a sine positional encoding of its pixel coordinates, a learned embedding of its pyramid
level and a learned embedding of its camera, so that \emph{every feature carries where and in which image it was
computed}.

\subsection{Prediction Branches}
\label{subsec:design-branches}

Every decoder layer has its own set of four prediction branches applied to its output embedding. The
classification branch is an \ac{mlp} of two hidden layers that outputs $N_{\mathrm{cls}}$ logits, and the
three regression branches are \acp{mlp} of the same depth that output six geometric offsets, the sine and
cosine of the heading, and the velocity $(v_x, v_y)$. Together they give the \textbf{box code}

\begin{equation}
    \vec{b} = \left( c_x,\, c_y,\, \log w,\, \log l,\, c_z,\, \log h,\, \sin\theta,\, \cos\theta,\,
    v_x,\, v_y \right).
    \label{eq:design-boxcode}
\end{equation}

The centre is predicted as an offset from the reference point in inverse-sigmoid space,
$\hat{\tilde{c}}_x = \sigma(\Delta c_x + r_x)$ and likewise for $y$ and $z$, and is then scaled back to
metres. The size (width, length, height as w, l, h) is predicted as an offset in logarithmic space and added to the size carried by the slot,
$\hat{\mathbf{z}} = \mathbf{z} + \Delta \mathbf{z}$. The class score of a slot is the maximum sigmoid
probability over the classes, $s_i = \max_k \sigma(\hat{y}_{i,k})$, taken from the \emph{last decoder layer}. Every
layer is supervised separately, so that each layer learns to produce a usable prediction.

The predictions of layer $l$ in frame $t$ are supervised by the detection loss

\begin{equation}
    \mathcal{L}_{\mathrm{det}}^{t,l} = \lambda_{\mathrm{cls}}\, \mathcal{L}_{\mathrm{cls}}^{t,l} +
    \lambda_{\mathrm{box}}\, \mathcal{L}_{\mathrm{box}}^{t,l},
    \label{eq:design-detloss}
\end{equation}

with the weights $\lambda_{\mathrm{cls}} = 2$ and $\lambda_{\mathrm{box}} = 0.25$. Which object a slot has to
predict, if any, is decided by the identity assignment of Algorithm~\ref{alg:design-assign}. With
$\mathcal{A}$ the set of slots that receive a target in this assignment and $\mathcal{P} \subseteq \mathcal{A}$
those whose target $g_i$ is an annotated object, the two terms are

\begin{equation}
    \begin{aligned}
        \mathcal{L}_{\mathrm{cls}}^{t,l} &= \frac{1}{\lvert \mathcal{A} \rvert} \sum_{i}
        \sum_{k=1}^{N_{\mathrm{cls}}} \operatorname{FL}\!\left( \sigma(\hat{y}_{i,k}),\, y_{i,k} \right), \\
        \mathcal{L}_{\mathrm{box}}^{t,l} &= \frac{1}{\lvert \mathcal{P} \rvert} \sum_{i \in \mathcal{P}}
        \sum_{j=1}^{10} \omega_j \left\lvert \hat{b}_{i,j} - b_{g_i,j} \right\rvert,
    \end{aligned}
    \label{eq:design-detterms}
\end{equation}

where the first sum runs over all slots of the frame, $y_{i,k} = 1$ if slot $i$ is assigned to an object of
class $k$ and $0$ otherwise, and $b_{g_i,j}$ is the $j$-th component of the box code of the assigned object.
The first eight components are weighted by $\omega_j = 1$ and the two velocity components by
$\omega_j = 0.2$. A slot whose object has left the scene belongs to $\mathcal{A}$ but not to $\mathcal{P}$, so it
is trained as background and receives no box target. The classification term uses the sigmoid focal loss
\cite{Lin2017Focal},

\begin{equation}
    \operatorname{FL}(p, y) =
    \begin{cases}
        -\alpha_{\mathrm{f}} \left( 1 - p \right)^{\gamma} \log p, & y = 1, \\[0.4em]
        -\left( 1 - \alpha_{\mathrm{f}} \right) p^{\gamma} \log \left( 1 - p \right), & y = 0,
    \end{cases}
    \label{eq:design-focal}
\end{equation}

with $\gamma = 2$ and $\alpha_{\mathrm{f}} = 0.25$. The factor $(1 - p)^{\gamma}$, or $p^{\gamma}$ for a
negative, lowers the loss of slots that are already classified correctly, so that the many empty slots that are
easy background do not outweigh the few slots that hold an object.

\subsection{Tracking}
\label{subsec:design-tracking}

Tracking decides which slot follows which object, keeps the embedding of a followed object informed about
its past, and moves its query into the next frame. It is realised differently in training and at inference.

\paragraph{Identity assignment in training.}
During training the identity of a slot is taken from the \ac{gt}, as in MOTR \cite{Zeng2022MOTR}. 
A slot that follows an object keeps it as long as the object is annotated, and only the empty slots compete for objects that no slot
follows yet. The competition is resolved by the Hungarian method \cite{Kuhn1955Hungarian}.

The procedure is given in
Algorithm~\ref{alg:design-assign}. It is run for \emph{each decoder layer separately} to supervise that layer, but
only the assignment made on the last layer is kept as the identity of the slot.

\begin{algorithm}[htbp]
    \caption{Identity assignment and supervision targets for one frame during training}
    \label{alg:design-assign}
    \begin{algorithmic}[1]
        \Require slots with identities $\mathrm{id}_i$ and predictions $(\hat{\mathbf{y}}_i, \hat{\vec{b}}_i)$,
        annotated objects $\mathcal{G}_t$ of frame $t$ with persistent identities
        \Ensure target $g_i$ of every slot, an object of $\mathcal{G}_t$ or none
        \For{every slot $i$ with $\mathrm{id}_i \geq 0$}
            \If{an object of $\mathcal{G}_t$ has identity $\mathrm{id}_i$}
                \State $g_i \gets$ that object \Comment{inherited, no matching needed}
            \Else
                \State $g_i \gets$ none \Comment{object left the scene, trained as background}
            \EndIf
        \EndFor
        \State $\mathcal{E} \gets \{ i : \mathrm{id}_i = -1 \}$ \Comment{empty slots}
        \State $\mathcal{U} \gets$ objects of $\mathcal{G}_t$ that no slot follows \Comment{new or reappearing}
        \State $\pi \gets \Call{Hungarian}{\mathcal{C}(\mathcal{E}, \mathcal{U})}$ \Comment{cost of Eq.~\eqref{eq:design-cost}}
        \For{every pair $(i, n) \in \pi$}
            \State $\mathrm{id}_i \gets$ identity of $n$, \quad $g_i \gets n$ \Comment{track birth}
        \EndFor
        \State classification loss on all slots, positive where $g_i$ is an object
        \State box loss on the slots whose $g_i$ is an object
    \end{algorithmic}
\end{algorithm}

\paragraph{Track birth and death at inference.}
At inference no annotation is available, and identities are governed by the score of each slot, as given in
Algorithm~\ref{alg:design-lifecycle}. An empty slot whose score reaches $\tau_{\mathrm{new}} = 0.4$ becomes a
new track. A track whose score falls below $\tau_{\mathrm{keep}} = 0.35$ accumulates a miss, and after
$m_{\max} = 5$ consecutive misses it is ended. Between the two thresholds a track neither resets nor
accumulates misses, so the gate has a hysteresis. Only tracks with a score of at least $\tau_{\mathrm{keep}}$
are reported in the output. These thresholds sit behind hard comparisons, so no gradient reaches them, and
they are inherited from ViP3D unchanged.

\begin{algorithm}[htbp]
    \caption{Track birth and death at inference}
    \label{alg:design-lifecycle}
    \begin{algorithmic}[1]
        \Require slots with scores $s_i$, identities $\mathrm{id}_i$ and miss counters $m_i$
        \For{every slot $i$}
            \If{$s_i \geq \tau_{\mathrm{new}}$}
                \State $m_i \gets 0$
            \EndIf
            \If{$\mathrm{id}_i = -1$ \textbf{and} $s_i \geq \tau_{\mathrm{new}}$}
                \State $\mathrm{id}_i \gets$ next unused identity \Comment{track birth}
            \ElsIf{$\mathrm{id}_i \geq 0$ \textbf{and} $s_i < \tau_{\mathrm{keep}}$}
                \State $m_i \gets m_i + 1$
                \If{$m_i \geq m_{\max}$}
                    \State $\mathrm{id}_i \gets -1$ \Comment{track death}
                \EndIf
            \EndIf
        \EndFor
    \end{algorithmic}
\end{algorithm}

%figure: Event (state) diagram of one slot at inference. States: empty, active, missing, ended. Transitions:
% empty to active when s >= 0.4 (birth), active to missing when s < 0.35, missing back to active when s >= 0.4,
% missing stays missing while 0.35 <= s < 0.4, missing to ended after 5 consecutive misses. A note that only
% slots with s >= 0.35 are reported.

\paragraph{Memory bank.}
Every track keeps its last $L$ stored embeddings in $\mathbf{B}_i$. Once a track has at least one stored
entry, its current output embedding $\mathbf{h}_i$ attends to them with multi-head attention, followed by a
residual connection, layer normalisation and an \ac{ffn}, so that the embedding of a track is refined by its
own history. A projection of the refined embedding is then pushed into the bank and the oldest entry dropped.
During training every frame is stored, while at inference a track stores an entry only every fourth frame,
so at test time the bank spans a longer stretch of the past than it does in training.

\paragraph{Query interaction module.}
The \ac{qim} updates the queries of the tracks that survive into the next frame. The active tracks attend to
each other through self-attention, with the sum of their positional half and output embedding as queries
and keys and their output embedding as values, followed by an \ac{ffn}. Two further \acp{ffn} map the result
to updates of the positional half and of the content half of the query, each added with a residual
connection and normalised. The query of a track thus carries what the decoder saw of it in this frame into
the next one. During training the module also drops each active track with probability 0.1 and adds, for
each active track with probability 0.3, one of the empty slots with the highest score, so that the model
learns to recover a lost track and to reject a false one.

\paragraph{Propagation into the next frame.}
Before a track is decoded in the next frame, its reference point is moved by its own predicted velocity and
by the motion of the ego vehicle. With $(\mathbf{R}_t, \mathbf{t}_t)$ the pose of the \ac{lidar} frame at
time $t$ in global coordinates and $\Delta t$ the time between the frames, the metric reference point
$\mathbf{x}_t$ becomes

\begin{equation}
    \mathbf{x}_{t+1} = \mathbf{R}_{t+1}^{\top} \left( \mathbf{R}_t \left( \mathbf{x}_t + \mathbf{v}_t\,
    \Delta t \right) + \mathbf{t}_t - \mathbf{t}_{t+1} \right),
    \label{eq:design-propagate}
\end{equation}

where $\mathbf{v}_t = (v_x, v_y, 0)$ is the velocity predicted by the last decoder layer. The result is
normalised and passed through $\sigma^{-1}$ as in Eq.~\eqref{eq:design-invsig}. When $\Delta t$ lies outside
$[0, 1]$\,s, as happens at the start of a new scene, only the ego transformation is applied.

\subsection{Trajectory Prediction}
\label{subsec:design-prediction}

Trajectory prediction follows VectorNet \cite{VectorNet2020} with the agent and lane interaction of LaneGCN
\cite{Liang2020LaneGCN}. The lanes and lane connectors of the \ac{hdmap} within 70\,m of the ego vehicle are
discretised at 1\,m, each becomes a polyline of vectors, and a polyline subgraph of self-attention layers and
max-pooling encodes each polyline into one vector.

The agents are not described by their past trajectories. Their representation is the output embedding of
their track query, which already holds what the decoder and the memory bank gathered about the agent, so that
nothing is lost at the interface between tracking and prediction. When camera features are available, this
embedding is first refined by one more projective camera cross-attention at the last box centre of the
agent. It then passes two \ac{mlp} layers, and the predictor reads its first 128 components, which is the
working width of the predictor. The lanes attend to the lanes and the agents, a global self-attention graph
relates all agents and lanes to each other, and a decoder outputs for every agent $M$ trajectories of $T_f$
positions in the coordinate frame of the agent, together with a probability $\pi_m$ for each of them.

The predictor is trained with a variety loss. For an agent $a$ whose last valid future position is at step
$\tau_a$, only the mode closest to the \ac{gt} at that step,

\begin{equation}
    m^{*} = \arg\min_{m} \left\lVert \hat{\mathbf{y}}^{a,m}_{\tau_a} - \mathbf{y}^{a}_{\tau_a} \right\rVert_2,
    \label{eq:design-bestmode}
\end{equation}

is regressed, and the probability of that mode is pushed up,

\begin{equation}
    \mathcal{L}_{\mathrm{pred}} = \frac{1}{A} \sum_{a=1}^{A} \left( \frac{1}{\lvert \mathcal{V}_a \rvert}
    \sum_{t \in \mathcal{V}_a} \operatorname{SmoothL1}\!\left( \hat{\mathbf{y}}^{a,m^{*}}_t -
    \mathbf{y}^{a}_t \right) - \log \pi^{a}_{m^{*}} \right),
    \label{eq:design-predloss}
\end{equation}

where $\mathcal{V}_a$ is the set of valid future steps of agent $a$ and $A$ the number of supervised agents.
Regressing only the closest mode lets the modes specialise to different futures instead of collapsing onto
their mean. Only tracks matched to an annotated agent with a valid future are supervised.


% ============================================================================
\section{Requirements}
\label{sec:design-requirements}
% ============================================================================

The requirements in this section apply to the parts that this thesis replaces or adds, and to how these parts
are joined with the components reused from ViP3D. They follow from the gaps and the starting points of
Section~\ref{sec:summary}.

\subsection{Functional Requirements}
\label{sec:design-fr}

These functional requirements state what the model has to do.

TODO: add FR0 for lidar modality
\begin{enumerate}[label=\textbf{FR\arabic*.}, ref=FR\arabic*, leftmargin=*]
    \item Every query slot that \emph{does not follow an object} is initialised from the current
    \ac{lidar} observation, with its reference point placed where an object is likely and its content read
    from the \ac{bev} features at that position.
    \item A query that already follows an object \emph{keeps its state from the previous frame}, and
    no new query is placed where such an object lies, so that the initialisation concerns newly appearing
    objects alone.
    \item The \emph{camera contributes through a fusion layer}  that refines queries already
    formed from the point cloud, and the model remains functional without that layer.
    \item Detection, tracking and trajectory prediction are \emph{produced by one network} from one
    set of queries in an \acl{e2e} manner, as in ViP3D.
\end{enumerate}

\subsection{Non-Functional Requirements}
\label{sec:design-nfr}

These non-functional requirements constrain how the model has to be implemented.

\begin{enumerate}[label=\textbf{NFR\arabic*.}, ref=NFR\arabic*, leftmargin=*]
    \item \textbf{Experimental control.} Every part of ViP3D \emph{outside the query path is kept unchanged}, so that a difference
    in the results can be attributed to the proposed change.
    \item \textbf{Architectural flexibility.} The decoder shall be built from layers that \emph{can be added or
    removed} without changing the rest of the model, so that the number of layers can be lowered below the six
    layers of ViP3D as far as the experiments show it to be sufficient.
    \item \textbf{Modularity and trainability.} The camera fusion layer shall be \emph{trainable independently} while 
    freezing the pre-trained \ac{lidar} feature extractor weights, ensuring the camera stream can be added on top of 
    a fixed geometric representation.
\end{enumerate}


% ============================================================================
\section{LiDAR Branches}
\label{sec:design-lidar}
% ============================================================================

The point cloud of the current frame is accumulated with the nine preceding sweeps, so that ten sweeps enter
the branch, and every point carries five features, its three coordinates, its intensity and its time offset
from the current sweep. The branch follows PointPillars \cite{Lang2019PointPillars} and the convolutional
backbone of SECOND \cite{Yan2018SECOND}, as introduced in Section~\ref{sec:rw-lidar}.

Space within $\pm 51.2$\,m around the vehicle is divided into pillars of $0.2 \times 0.2$\,m that span the
whole vertical range, and a pillar holds at most 20 points. A PointNet-style encoder \cite{Qi2017PointNet}
turns the points of each non-empty pillar into a vector of 64 features, and the vectors are scattered back to
the positions of their pillars, which gives a pseudo-image of $512 \times 512$ cells. A two-dimensional
\ac{cnn} with two strided stages reduces it to $256 \times 256$ cells with 128 channels and to $128 \times 128$
cells with 256 channels, and a neck upsamples both to $256 \times 256$ and concatenates them. The result is
the \ac{bev} feature map

\begin{equation}
    \mathbf{F} \in \mathbb{R}^{C \times H \times W}, \qquad C = 384,\quad H = W = 256,\quad s = 0.4\,\mathrm{m}.
    \label{eq:design-bev}
\end{equation}

The map serves twice. The heatmap head predicts object centres from it, and the first decoder layer reads
features from it at the reference points of the queries.

%figure: Block scheme of the LiDAR branch. Ten accumulated sweeps (N points x 5 features) -> pillar grid
% 512 x 512 at 0.2 m -> PointNet pillar encoder (5 -> 64) -> scatter to pseudo-image 64 x 512 x 512 -> 2D CNN
% stage 1 (stride 2, 128 x 256 x 256) and stage 2 (stride 2, 256 x 128 x 128) -> neck upsampling both to
% 256 x 256 -> concatenation, BEV map 384 x 256 x 256 at 0.4 m, which feeds both the heatmap head and decoder
% layer 0.


% ============================================================================
\section{Query Initialisation from the LiDAR Heatmap}
\label{sec:design-init}
% ============================================================================

This section describes the first contribution of the thesis, which replaces the learned embedding of ViP3D
by queries derived from the point cloud, as proposed in Starting point~1. The principle was introduced in
Section~\ref{sec:rw-lidar}, and what follows is its design inside a tracking model.

\subsection{Class-Wise Heatmap Head}
\label{subsec:design-heatmap}

The heatmap head predicts, for every class and every \ac{bev} cell, the probability that the centre of an
object of that class lies in the cell. A shared $3 \times 3$ convolution reduces the \ac{bev} map to 64
channels, followed by batch normalisation and a rectified linear unit. Several small heads then each predict
a group of classes, with the classes grouped as in CenterPoint \cite{Yin2021CenterPoint} so that similar
shapes share a head. Each head consists of a $3 \times 3$ convolution with normalisation and activation and a
final $3 \times 3$ convolution with one output channel per class. The outputs are reordered into the class
order of the model, which gives the logits $\hat{\mathbf{Y}} \in \mathbb{R}^{N_{\mathrm{cls}} \times H
\times W}$. The bias of the final convolutions is initialised so that every cell starts with a probability
of about 0.1, which keeps the loss stable in the first iterations when almost every cell is background.

The target heatmap is built from the annotated boxes. An object $n$ of class $k_n$ with centre $(x_n, y_n)$
and footprint $w_n \times l_n$ falls into the cell $(i_n, j_n)$ with $j_n = \lfloor (x_n - x_{\min}) / s
\rfloor$ and $i_n = \lfloor (y_n - y_{\min}) / s \rfloor$, and is drawn as a Gaussian whose radius follows
its size,

\begin{equation}
    \begin{aligned}
        Y_{k,i,j} &= \max_{n:\, k_n = k} \exp\!\left( -\frac{(i - i_n)^2 + (j - j_n)^2}{2 \sigma_n^2} \right), \\
        \sigma_n &= \frac{2\rho_n + 1}{6}, \qquad
        \rho_n = \max\!\left( 1, \left\lceil \frac{\sqrt{w_n^2 + l_n^2}}{2s} \right\rceil \right),
    \end{aligned}
    \label{eq:design-target}
\end{equation}

evaluated within a window of $2\rho_n + 1$ cells around the centre. The radius $\rho_n$ is half the diagonal
of the footprint in cells, so a bus spreads over more cells than a pedestrian, and the maximum keeps the
peak of every object at exactly 1 where objects of the same class overlap. The head is trained with the
Gaussian focal loss of CornerNet \cite{Law2018CornerNet},

\begin{equation}
    \mathcal{L}_{\mathrm{hm}} = -\frac{1}{N_{\mathrm{pos}}} \sum_{k,i,j}
    \begin{cases}
        \left( 1 - \hat{P}_{k,i,j} \right)^{\alpha} \log \hat{P}_{k,i,j}, & Y_{k,i,j} = 1, \\[0.4em]
        \left( 1 - Y_{k,i,j} \right)^{\beta} \hat{P}_{k,i,j}^{\,\alpha} \log\!\left( 1 - \hat{P}_{k,i,j} \right),
        & \text{otherwise},
    \end{cases}
    \label{eq:design-hmloss}
\end{equation}

where $\hat{\mathbf{P}} = \sigma(\hat{\mathbf{Y}})$, $N_{\mathrm{pos}}$ is the number of cells with
$Y_{k,i,j} = 1$, and $\alpha = 2$ and $\beta = 4$. The factor $(1 - Y)^{\beta}$ lowers the penalty on cells
close to a centre, so the head is not punished for a peak that lands one cell off. Normalising by
$N_{\mathrm{pos}}$ makes the loss an average over objects instead of a sum over the whole grid.

%figure: Visualisation of the heatmap on one frame: the BEV feature map, the predicted heatmap (maximum over
% classes) and the target Gaussians side by side, with the selected peaks marked and the ground-truth boxes
% drawn. The figure presentation/figures/lidar-bev-pred.png can be reused.

\subsection{Selecting Peaks and Building Queries}
\label{subsec:design-select}

The empty slots of a frame are filled by Algorithm~\ref{alg:design-init}. The heatmap probabilities are
first detached from the computation graph, since the selection of a cell is a discrete choice through which
no gradient can pass. Cells around every active track are then set to zero within a window of $5 \times 5$
cells, which is $2 \times 2$\,m, so that a new query cannot be born on an object that a track already
follows. A $3 \times 3$ max-pooling keeps only the local maxima of every class, which acts as a cheap
non-maximum suppression. It is skipped for bicycles and pedestrians, whose footprints are smaller than the
pooling window and which often stand closer to each other than one window, so that suppression would remove
real objects. The strongest cells are finally selected jointly over all classes and all positions, so that
the classes compete for the slots according to their scores.

A selected cell $(i, j)$ of class $k$ becomes a query in two steps. Its reference point is the metric centre
of the cell,

\begin{equation}
    x = x_{\min} + \left( j + \tfrac{1}{2} \right) s, \qquad
    y = y_{\min} + \left( i + \tfrac{1}{2} \right) s, \qquad
    z = \tfrac{1}{2} \left( z_{\min} + z_{\max} \right),
    \label{eq:design-cell2metric}
\end{equation}

normalised and passed through $\sigma^{-1}$ as in Eq.~\eqref{eq:design-invsig}. The height is set to the
middle of the range, since the \ac{bev} map carries no vertical position, and is left for the decoder to
refine. Its content is the \ac{bev} feature of the cell, projected to the embedding width, plus a learned
embedding of the class,

\begin{equation}
    \mathbf{c} = \mathbf{W}_{\mathrm{bev}}\, \mathbf{F}_{:,i,j} + \mathbf{e}_k, \qquad
    \mathbf{p} = \mathbf{c},
    \label{eq:design-content}
\end{equation}

where $\mathbf{W}_{\mathrm{bev}} \in \mathbb{R}^{D \times C}$ and $\mathbf{e}_k \in \mathbb{R}^{D}$. The
positional half is set equal to the content half at birth, since the content already encodes where the cell
is, and the two halves diverge as soon as the query interaction module updates them.

\begin{algorithm}[htbp]
    \caption{Initialisation of the empty slots from the \ac{lidar} heatmap}
    \label{alg:design-init}
    \begin{algorithmic}[1]
        \Require heatmap logits $\hat{\mathbf{Y}} \in \mathbb{R}^{N_{\mathrm{cls}} \times H \times W}$,
        \ac{bev} map $\mathbf{F}$, slots $\mathcal{S}$
        \Ensure the empty slots of $\mathcal{S}$ receive a reference point and a query
        \Function{InitQueries}{$\hat{\mathbf{Y}}, \mathbf{F}, \mathcal{S}$}
            \State $\mathbf{P} \gets \operatorname{stopgrad}\big( \sigma(\hat{\mathbf{Y}}) \big)$
            \For{every slot $a \in \mathcal{S}$ with $\mathrm{id}_a \geq 0$}
                \State $(i_a, j_a) \gets$ \ac{bev} cell of the reference point $\mathbf{r}_a$
                \State $\mathbf{P}_{:,\, i_a-2 : i_a+2,\; j_a-2 : j_a+2} \gets 0$ \Comment{no birth on a tracked object}
            \EndFor
            \State $\mathbf{P}^{\max} \gets \operatorname{maxpool}_{3 \times 3}(\mathbf{P})$
            \For{every class $k$ except bicycle and pedestrian}
                \State $\mathbf{P}_k \gets \mathbf{P}_k \odot \mathbf{1}\!\left[ \mathbf{P}_k = \mathbf{P}^{\max}_k \right]$
                \Comment{keep local maxima only}
            \EndFor
            \State $\mathcal{E} \gets \{ i \in \mathcal{S} : \mathrm{id}_i = -1 \}$ \Comment{empty slots}
            \State $\{(k_n, i_n, j_n)\}_{n=1}^{\lvert \mathcal{E} \rvert} \gets$ the $\lvert \mathcal{E} \rvert$
            largest entries of $\mathbf{P}$ over all classes and cells
            \For{$n = 1, \dots, \lvert \mathcal{E} \rvert$}
                \State $(x, y, z) \gets$ metric centre of cell $(i_n, j_n)$ \Comment{Eq.~\eqref{eq:design-cell2metric}}
                \State $\mathbf{r}_{\mathcal{E}_n} \gets \sigma^{-1}\big( \operatorname{normalise}(x, y, z) \big)$
                \State $\mathbf{c} \gets \mathbf{W}_{\mathrm{bev}}\, \mathbf{F}_{:,i_n,j_n} + \mathbf{e}_{k_n}$
                \State $(\mathbf{p}_{\mathcal{E}_n}, \mathbf{c}_{\mathcal{E}_n}) \gets (\mathbf{c}, \mathbf{c})$
            \EndFor
        \EndFunction
    \end{algorithmic}
\end{algorithm}

\subsection{Integration with the Tracker}
\label{subsec:design-integration}

The initialisation is inserted into ViP3D at the single place where an empty slot used to receive its learned
values, and it touches no other part of the tracking. Four properties of the integration follow from that
placement.

\begin{itemize}
    \item \textbf{Only empty slots are filled.} A slot with an identity keeps the query and the reference
    point it carried from the previous frame, so a followed object is never re-initialised and keeps its
    history. The suppression around active tracks complements this by keeping new queries off the cells
    those tracks occupy.
    \item \textbf{Birth is still decided by the tracker.} Placing a query on a peak does not create a track.
    In training the slot must still win the matching of Algorithm~\ref{alg:design-assign}, and at inference
    its score must still reach $\tau_{\mathrm{new}}$ in Algorithm~\ref{alg:design-lifecycle}. A weak peak
    therefore produces a query that the decoder can reject.
    \item \textbf{Gradients reach the initialisation only through the content.} The choice of cells is
    discrete, so the heatmap head learns only from its own loss of Eq.~\eqref{eq:design-hmloss}. The
    detection loss reaches the \ac{bev} features, the projection $\mathbf{W}_{\mathrm{bev}}$ and the class
    embeddings through the content of Eq.~\eqref{eq:design-content}, which teaches them to describe an object
    in a way the decoder can use.
    \item \textbf{The initial box size is a learned constant.} The size of a new slot is not derived from
    its class, so every new query starts from the same learned size, which the first decoder layer refines
    before the second layer needs it.
\end{itemize}

Slots that the query interaction module adds as false positives during training keep the identity $-1$, so
in the next frame they are re-initialised from the heatmap like any other empty slot. The complete
processing of one frame at inference is summarised in Algorithm~\ref{alg:design-frame}.

\begin{algorithm}[htbp]
    \caption{Processing of one frame by \modelname at inference}
    \label{alg:design-frame}
    \begin{algorithmic}[1]
        \Require point cloud $\mathcal{P}_t$, images $\{I^m_t\}_{m=1}^{N_{\mathrm{cam}}}$, tracks
        $\mathcal{T}_{t-1}$, ego poses of frames $t-1$ and $t$
        \Ensure tracks $\mathcal{T}_t$ with boxes, classes and identities
        \State move the reference points of $\mathcal{T}_{t-1}$ into frame $t$ \Comment{Eq.~\eqref{eq:design-propagate}}
        \State $\mathcal{S} \gets \mathcal{T}_{t-1} \cup \{ N_q \text{ empty slots} \}$
        \State $\mathbf{F} \gets \Call{LidarBranch}{\mathcal{P}_t}$, \quad
        $\{\mathbf{F}^m\} \gets \Call{CameraBranch}{\{I^m_t\}}$
        \State $\hat{\mathbf{Y}} \gets \Call{HeatmapHead}{\mathbf{F}}$
        \State \Call{InitQueries}{$\hat{\mathbf{Y}}, \mathbf{F}, \mathcal{S}$} \Comment{Algorithm~\ref{alg:design-init}}
        \State decode all slots of $\mathcal{S}$ with layer 0 and layer 1 \Comment{Section~\ref{sec:design-decoder}}
        \State update scores and identities \Comment{Algorithm~\ref{alg:design-lifecycle}}
        \State refine the embeddings of the slots with the memory bank
        \State $\mathcal{T}_t \gets$ the slots with $\mathrm{id} \geq 0$, their queries updated by the \ac{qim}
        \State \Return $\mathcal{T}_t$
    \end{algorithmic}
\end{algorithm}


% ============================================================================
\section{Two-Layer Fusion Decoder}
\label{sec:design-decoder}
% ============================================================================

This section describes the second contribution, the shallow decoder of Starting point~2. It replaces the six
camera layers of ViP3D with two layers of different kinds. The first reads the \ac{bev} map and the second
attends to the camera features, which follows the decoder of TransFusion \cite{Bai2022TransFusion} described
in Section~\ref{sec:rw-fusion}. Both layers share the structure of a transformer decoder layer, namely
self-attention among all queries, cross-attention to a sensor, and an \ac{ffn} with a hidden width of 512,
each followed by a residual connection and layer normalisation, with eight attention heads throughout.

%figure: Internal structure of the two-layer decoder. Input: queries (positional and content halves),
% reference points and box sizes. Layer 0: self-attention (8 heads) -> add & norm -> BEV read at the reference
% point -> add & norm -> FFN 256-512-256 -> add & norm -> prediction branches 0 (cls 7, reg 6, dir 2, vel 2) and
% refinement of reference point and size. Layer 1: self-attention -> add & norm -> SMCA over the camera features
% -> add & norm -> FFN -> add & norm, with a fallback arrow from the layer 0 output for queries visible in no
% camera -> prediction branches 1 and refinement. Can be drawn from the hand-written decoder sketch.

\subsection{Layer 0: Reading the BEV Map}
\label{subsec:design-layer0}

The cross-attention of the first layer has the form of deformable attention \cite{Zhu2021DeformableDETR}
with $P = 4$ sampling points per query and learned weights over them. In the configuration used, the learned
sampling offsets are disabled, so all four points coincide with the reference point of the query and the
layer reduces to a bilinear read of the \ac{bev} map at that point,

\begin{equation}
    \mathbf{o}_i = \mathbf{W}_o\, \mathbf{W}_{\mathrm{in}}\, \mathbf{F}\!\left( \tilde{r}_{i,x},
    \tilde{r}_{i,y} \right),
    \label{eq:design-layer0}
\end{equation}

where $\mathbf{F}(\cdot)$ denotes bilinear interpolation of the map at the normalised position,
$\mathbf{W}_{\mathrm{in}} \in \mathbb{R}^{D \times C}$ projects the \ac{bev} channels to the embedding width
and $\mathbf{W}_o$ is an output projection. The output is added to the query as a residual. The reasoning
behind this simplification is that a new query is born on a heatmap peak and a track has just been moved
onto its predicted position, so in both cases the reference point already lies on the object and the
feature at that point is the evidence the query needs. Reading a fixed number of points also makes the cost
of the layer independent of the size of the \ac{bev} grid.

\subsection{Layer 1: Spatially Modulated Camera Attention}
\label{subsec:design-layer1}

The second layer reintroduces the camera through \ac{smca} \cite{Gao2021SMCA}, as proposed in Starting
point~2. Unlike the projective cross-attention of ViP3D, which reads one point per camera, it attends to every
position of a camera feature map, and a Gaussian centred on the projection of the reference point decides
how much each position may contribute. It reads only the finest pyramid level, at stride 8, whose map has
$h \times w$ positions per camera.

For each camera $m$, the reference point is projected into the image, which gives its position
$(u_m, v_m)$ on the feature map and a flag $\nu_m$ that holds when the point lies in front of the camera and
inside the image. The width of the Gaussian follows the size of the box in the image. The centre of the box
and the four points at half its width and half its length along the axes are projected, the extent they
cover on the feature map is measured, and

\begin{equation}
    \varsigma_m = \operatorname{clip}\!\left( \beta \cdot \tfrac{1}{2} \sqrt{\Delta u_m^2 + \Delta v_m^2},\;
    1,\; 50 \right), \qquad
    G_m(u, v) = \exp\!\left( -\frac{(u - u_m)^2 + (v - v_m)^2}{\varsigma_m^2} \right),
    \label{eq:design-sigma}
\end{equation}

where $\Delta u_m$ and $\Delta v_m$ are the horizontal and vertical extent in feature-map positions,
$\beta = 1$ is a scale, and only points in front of the camera are used to measure the extent. A nearby car
covers many positions and receives a wide Gaussian, a distant one a narrow Gaussian. The attention of each
head adds the logarithm of the Gaussian to the logits before the softmax,

\begin{equation}
    \mathbf{A}_m = \operatorname{softmax}\!\left( \frac{\mathbf{Q} \mathbf{K}_m^{\top}}{\sqrt{d_h}} +
    \log \mathbf{G}_m \right), \qquad
    \mathbf{o}_m = \mathbf{A}_m \mathbf{V}_m,
    \label{eq:design-smca}
\end{equation}

where $\mathbf{Q}$ is projected from the queries, $\mathbf{K}_m$ and $\mathbf{V}_m$ from the features of
camera $m$, and $d_h = D / 8$. Since $\log G_m = 0$ at the projected centre and falls quadratically away from
it, a position far from the object is excluded however similar its features are. The outputs are averaged
over the cameras that see the query,

\begin{equation}
    \mathbf{o} = \frac{\sum_{m} \nu_m\, \mathbf{o}_m}{\max\!\left( 1, \sum_{m} \nu_m \right)},
    \label{eq:design-camavg}
\end{equation}

projected, and added to the query together with an encoding of the three-dimensional reference point
produced by a small \ac{mlp}. A query that no camera sees receives nothing from Eq.~\eqref{eq:design-camavg},
and for such queries the whole output of the layer is replaced by the output of layer 0. Objects outside the
field of view of every camera, or close to the vehicle below the lower image edges, are therefore described
by the \ac{lidar} alone rather than by an update computed from no evidence. The computation is summarised in
Algorithm~\ref{alg:design-smca}.

\begin{algorithm}[htbp]
    \caption{Spatially modulated cross-attention of decoder layer 1}
    \label{alg:design-smca}
    \begin{algorithmic}[1]
        \Require queries $\mathbf{X}$, reference points $\mathbf{r}$, box sizes $\mathbf{z}$,
        camera features $\mathbf{F}^1, \dots, \mathbf{F}^{N_{\mathrm{cam}}}$
        \Ensure updated queries and the set of queries seen by no camera
        \Function{SMCA}{$\mathbf{X}, \mathbf{r}, \mathbf{z}, \{\mathbf{F}^m\}$}
            \State $\mathbf{o} \gets \mathbf{0}$, \quad $\mathbf{n} \gets \mathbf{0}$
            \State $\mathbf{Q} \gets \mathbf{W}_Q \mathbf{X}$ split into 8 heads
            \For{$m = 1, \dots, N_{\mathrm{cam}}$}
                \State $(u_m, v_m, \nu_m) \gets$ projection of the reference points into camera $m$
                \State $\varsigma_m \gets$ width from the projected box extent \Comment{Eq.~\eqref{eq:design-sigma}}
                \State $\mathbf{G}_m \gets$ Gaussian over the $h \times w$ positions
                \State $\mathbf{K}_m, \mathbf{V}_m \gets \mathbf{W}_K \mathbf{F}^m,\ \mathbf{W}_V \mathbf{F}^m$
                \State $\mathbf{A}_m \gets \operatorname{softmax}\big( \mathbf{Q}\mathbf{K}_m^{\top} / \sqrt{d_h} + \log \mathbf{G}_m \big)$
                \State $\mathbf{o} \gets \mathbf{o} + \nu_m\, \mathbf{A}_m \mathbf{V}_m$, \quad $\mathbf{n} \gets \mathbf{n} + \nu_m$
            \EndFor
            \State $\mathbf{o} \gets \mathbf{W}_o\, \mathbf{o} / \max(1, \mathbf{n})$
            \State \Return $\mathbf{X} + \mathbf{o} + \operatorname{MLP}\big( \sigma^{-1}(\tilde{\mathbf{r}}) \big)$, \quad $\{ i : n_i = 0 \}$
        \EndFunction
    \end{algorithmic}
\end{algorithm}

%figure: SMCA on one frame: the six camera images with the projected reference points of the top 20 queries,
% the Gaussian of each drawn as an ellipse or heat overlay, and next to it the resulting attention maps. Shows that
% a near object gets a wide Gaussian and a far one a narrow Gaussian. The figure
% presentation/figures/smca-attn-top20.png can be reused.

\subsection{Reference Refinement}
\label{subsec:design-refine}

After each decoder layer $l$, the regression branch of that layer predicts an offset of the reference point
and of the box size, and both are updated before the next layer,

\begin{equation}
    \tilde{\mathbf{r}}^{(l+1)} = \operatorname{stopgrad}\!\left( \sigma\!\left( \sigma^{-1}\big(
    \tilde{\mathbf{r}}^{(l)} \big) + \Delta \mathbf{r}^{(l)} \right) \right), \qquad
    \mathbf{z}^{(l+1)} = \operatorname{stopgrad}\!\left( \mathbf{z}^{(l)} + \Delta \mathbf{z}^{(l)} \right),
    \label{eq:design-refine}
\end{equation}

where $\tilde{\mathbf{r}}^{(l)} \in [0, 1]^3$ is the normalised reference point entering layer $l$, so the
offset is added in inverse-sigmoid space and the result is mapped straight back into the range.

The stop-gradient keeps each layer responsible for its own prediction. A layer is supervised on the
correction it makes, and its error is not propagated back into the layer that supplied its starting point,
which stabilises training in the same way as in Deformable DETR. The refinement of layer 0 matters
particularly here, since layer 1 uses the refined reference point to place its Gaussian and the refined size
to set its width. The reference point produced by the last layer is the one that is propagated into the next
frame.

\subsection{Comparison with the Decoder of ViP3D}
\label{subsec:design-compare}

Table~\ref{tab:design-compare} sets the two decoders side by side. The reduction from six layers to two is
possible because the queries no longer start from constants, and the work of moving a query onto its object
is done before the decoder by the heatmap. The camera layer is, however, more expensive than a layer of ViP3D,
since it attends densely to whole feature maps instead of sampling one point per camera. The saving in depth
and the cost of the dense attention are therefore not independent, and Section~\ref{sec:design-complexity}
estimates both.

\begin{table}[htbp]
    \centering
    \caption{The decoder of ViP3D compared with the decoder of \modelname.}
    \label{tab:design-compare}
    \begin{tabular}{>{\raggedright\arraybackslash}p{0.32\textwidth}
                    >{\raggedright\arraybackslash}p{0.29\textwidth}
                    >{\raggedright\arraybackslash}p{0.29\textwidth}}
        \toprule
        & ViP3D & \modelname \\
        \midrule
        Initialisation of empty slots & learned embedding, the same in every frame & strongest heatmap peaks
        of the current frame \\
        Supervision before decoding & none & heatmap loss \\
        Number of layers & 6 & 2 \\
        \ac{lidar} & not used & read at the reference point in layer 0 \\
        Camera & projective sampling, one point per camera and level & dense attention over one level,
        modulated by a Gaussian \\
        Query seen by no camera & sampled features masked & output of layer 0 kept \\
        \bottomrule
    \end{tabular}
\end{table}


% ============================================================================
\section{Camera Visibility for Partial Camera Rigs}
\label{sec:design-visibility}
% ============================================================================

The six cameras of nuScenes cover the full circle around the vehicle. When the model is configured with
fewer cameras, objects behind the vehicle remain annotated but leave no trace in any image, and a camera
model would be trained to detect what it cannot see. Two masks prevent this, and both are active only when
fewer than six cameras are used. In training, an annotated box whose centre projects into no camera keeps
its place in the matching, so that the number of objects stays consistent, but its classification and box
losses are set to zero. In the selection of Algorithm~\ref{alg:design-init}, every \ac{bev} cell is projected
into the cameras, and cells that no camera sees are set to zero before the peaks are taken, so that queries
are born only where a camera can confirm them. With six cameras both masks are left out, since they would
otherwise also remove objects close to the vehicle that the \ac{lidar} sees perfectly well.


% ============================================================================
\section{Training}
\label{sec:design-training}
% ============================================================================

\subsection{Loss}
\label{subsec:design-loss}

The model is trained on clips of $T$ consecutive frames, so that identities can be carried from one frame to
the next and the tracking modules receive a learning signal. The loss of a clip is

\begin{equation}
    \mathcal{L} = \sum_{t=1}^{T} \left( \sum_{l=0}^{1} \mathcal{L}_{\mathrm{det}}^{t,l} +
    \lambda_{\mathrm{hm}}\, \mathcal{L}_{\mathrm{hm}}^{t} \right) + \mathcal{L}_{\mathrm{pred}},
    \label{eq:design-loss}
\end{equation}

where $\mathcal{L}_{\mathrm{det}}^{t,l}$ is the detection loss of Eq.~\eqref{eq:design-detloss} on the
predictions of layer $l$ in frame $t$, $\mathcal{L}_{\mathrm{hm}}^{t}$ is the heatmap loss of
Eq.~\eqref{eq:design-hmloss} with $\lambda_{\mathrm{hm}} = 1$, and $\mathcal{L}_{\mathrm{pred}}$ is the
prediction loss of Eq.~\eqref{eq:design-predloss}. Each detection and heatmap term is \emph{normalised} within its own
frame, the detection terms by the number of assigned slots and the heatmap term by the number of peaks, so that
a crowded frame does not outweigh a sparse frame.

\subsection{Two-Stage Schedule}
\label{subsec:design-schedule}

Training follows the two-stage recipe of TransFusion \cite{Bai2022TransFusion}, which satisfies NFR3.

\begin{itemize}
    \item \textbf{Stage 1, \ac{lidar} only.} The camera input is absent. Layer 1 is bypassed and its outputs
    are copies of those of layer 0, so the second set of prediction branches is already trained on a
    meaningful representation when the camera layer is introduced. Everything else is trained, namely the
    \ac{lidar} branch, the heatmap head, layer 0, both sets of prediction branches, the tracking modules and
    the predictor.
    \item \textbf{Stage 2, camera fusion.} Starting from the end of stage 1, the \ac{lidar} feature extractor
    is frozen and the camera layer is trained from its initialisation. The image backbone, pretrained for
    detection, is frozen as well, and the image neck is fine-tuned at a tenth of the learning rate. The
    heatmap head is fine-tuned at a tenth of the learning rate, while layer 0, the prediction branches, the
    tracking modules and the predictor continue to train at the full rate.
\end{itemize}

Freezing the \ac{lidar} feature extractor means that the camera cannot degrade the geometric representation
that the queries are built from. Fusion becomes a refinement of hypotheses that the point cloud has already
formed, as intended in Section~\ref{sec:rw-positioning}. Algorithm~\ref{alg:design-train} summarises the
processing of one training clip.

%figure: Diagram of the two training stages. Stage 1: LiDAR branch, heatmap head, layer 0, branches, tracking and
% predictor all trainable (one colour), camera path absent, layer 1 bypassed. Stage 2: LiDAR feature extractor and
% image backbone frozen (grey, padlock), image neck and heatmap head at 0.1 of the learning rate (light colour),
% SMCA layer newly trained (highlighted), layer 0, branches, tracking and predictor trainable.

\begin{algorithm}[htbp]
    \caption{Training on one clip}
    \label{alg:design-train}
    \begin{algorithmic}[1]
        \Require clip of $T$ frames with sensor data, annotated boxes with identities and future
        trajectories, stage $\in \{1, 2\}$
        \Ensure loss $\mathcal{L}$ of the clip
        \State $\mathcal{S} \gets N_q$ empty slots
        \For{$t = 1, \dots, T$}
            \State extract $\mathbf{F}$, and the camera features in stage 2
            \State $\hat{\mathbf{Y}} \gets \Call{HeatmapHead}{\mathbf{F}}$, \quad add $\lambda_{\mathrm{hm}} \mathcal{L}^{t}_{\mathrm{hm}}$ \Comment{Eq.~\eqref{eq:design-hmloss}}
            \State \Call{InitQueries}{$\hat{\mathbf{Y}}, \mathbf{F}, \mathcal{S}$}
            \State decode with layer 0, and with layer 1 in stage 2 or a copy of layer 0 in stage 1
            \For{$l = 0, 1$}
                \State assign targets to the predictions of layer $l$ \Comment{Algorithm~\ref{alg:design-assign}}
                \State add $\mathcal{L}_{\mathrm{det}}^{t,l}$ \Comment{Eq.~\eqref{eq:design-detloss}}
            \EndFor
            \State keep the identities assigned on layer 1
            \State memory bank, then \ac{qim} with random drop and false-positive insertion
            \If{$t < T$}
                \State move the reference points into frame $t+1$ \Comment{Eq.~\eqref{eq:design-propagate}}
            \EndIf
            \State $\mathcal{S} \gets$ surviving tracks $\cup \{ N_q \text{ empty slots} \}$
        \EndFor
        \State predict trajectories of the tracks matched to annotated agents, add $\mathcal{L}_{\mathrm{pred}}$ \Comment{Eq.~\eqref{eq:design-predloss}}
        \State \Return $\mathcal{L}$
    \end{algorithmic}
\end{algorithm}


% ============================================================================
\section{Complexity Analysis}
\label{sec:design-complexity}
% ============================================================================

Table~\ref{tab:design-complexity} estimates the time and memory complexity of each stage for one frame. Time
is given as the number of multiply and accumulate operations, and the orders of magnitude are computed from
the layer shapes of the configuration of Table~\ref{tab:design-notation}, with $Q \approx 300$. They estimate
the arithmetic of each stage and are not measured run times, which Chapter~\ref{chapter:experiments-and-evaluation}
reports.

\begin{table}[htbp]
    \centering
    \caption{Estimated time and memory complexity of one frame. $N_{\mathrm{pts}}$ is the number of points,
    $H_p W_p$ the pillar grid, $k$ the convolution kernel, $hw$ the positions of one camera feature map,
    $N_a$ the active tracks, $n$ the size of the assignment problem, and $A$ and $M_{\mathrm{l}}$ the agents
    and lanes of the predictor. Orders of magnitude refer to the configuration used.}
    \label{tab:design-complexity}
    \small
    \begin{tabular}{>{\raggedright\arraybackslash}p{0.25\textwidth}
                    >{\raggedright\arraybackslash}p{0.27\textwidth}
                    >{\raggedright\arraybackslash}p{0.17\textwidth}
                    >{\raggedright\arraybackslash}p{0.19\textwidth}}
        \toprule
        Stage & Time & Order of time & Memory \\
        \midrule
        Pillar encoding, scatter & $O(N_{\mathrm{pts}})$ & $10^{7}$ & $O(H_p W_p)$ \\
        \ac{lidar} \ac{cnn} and neck & $O(H W C^2 k^2)$ & $10^{11}$ & $O(H W C)$ \\
        Camera backbone and \ac{fpn} & $O(N_{\mathrm{cam}}\, H_{\mathrm{img}} W_{\mathrm{img}})$ & $10^{11}$ to $10^{12}$
        & $O(N_{\mathrm{cam}}\, hw D)$ \\
        Heatmap head & $O(H W C k^2)$ & $10^{9}$ & $O(N_{\mathrm{cls}} H W)$ \\
        Query selection & $O(N_{\mathrm{cls}} H W + N_a)$ & $10^{6}$ & $O(N_{\mathrm{cls}} H W)$ \\
        Self-attention, per layer & $O(Q^2 D)$ & $10^{7}$ & $O(Q^2)$ \\
        Layer 0, \ac{bev} read & $O(Q P C D)$ & $10^{8}$ & $O(Q P C)$ \\
        Layer 1, \ac{smca} & $O(N_{\mathrm{cam}}(Q\, hw D + hw D^2))$ & $10^{10}$ & $O(Q\, N_{\mathrm{cam}}\, hw)$ \\
        Prediction branches & $O(Q D^2)$ & $10^{8}$ & $O(Q D)$ \\
        Memory bank & $O(N_a L D)$ & $10^{6}$ & $O(N_a L D)$ \\
        Query interaction & $O(N_a^2 D)$ & $10^{6}$ & $O(N_a^2)$ \\
        Hungarian matching, training only & $O(n^3)$ & $10^{7}$ & $O(n^2)$ \\
        Trajectory prediction & $O((A + M_{\mathrm{l}})^2 D)$ & $10^{7}$ & $O((A + M_{\mathrm{l}})^2)$ \\
        \bottomrule
    \end{tabular}
\end{table}

Three observations follow from the table.

\paragraph{The query path is cheap.} The selection of queries from the heatmap is linear in the size of the
grid and costs no more than a pooling and a partial sort, and the read of layer 0 is linear in the number of
queries and independent of the grid. The contribution of this thesis to the query path therefore adds almost
nothing to the cost of a frame. The heatmap head costs more, but it is a few convolutions over a map that the
backbone has already produced.

\paragraph{The feature extractors dominate the arithmetic.} The two convolutional backbones account for the
largest share of the operations. The camera backbone is the larger of the two, since it processes six images,
and it is absent in stage 1, which makes the \ac{lidar}-only model considerably cheaper than the fused one
independently of the decoder.

\paragraph{The camera layer dominates the decoder and its memory.} Removing four layers saves little, since a
layer of ViP3D samples only a few points per query. The \ac{smca} layer instead forms, for each camera, an
attention matrix of $8 \times Q \times hw$ entries, about $2 \times 10^{7}$ values, and a Gaussian mask of
$Q \times N_{\mathrm{cam}} \times hw$ entries, about $1.5 \times 10^{7}$ values. These are the largest
intermediate tensors of the decoder, and since the layer moves much data for comparatively little arithmetic,
its run time is limited by memory traffic rather than by the operation count of the table. The shallower
decoder is therefore not cheaper by construction, and whether it is faster than the six layers of ViP3D is a
question for measurement.

\paragraph{Communication.} At inference the model runs on a single device and exchanges no data. Training is
data-parallel, so each step exchanges the gradients of the trainable parameters among the devices, which is
$O(\lvert \theta \rvert)$ per step and fewer parameters in stage 2, where the feature extractors are frozen.


% ============================================================================
\section{Summary}
\label{sec:design-summary}
% ============================================================================

\modelname keeps the agent queries, the tracking modules and the trajectory predictor of ViP3D, and changes
how a query begins and how it is refined. Empty slots are no longer filled from a learned constant but from
the strongest peaks of a class-wise heatmap predicted from the \ac{lidar} \ac{bev} map, with suppression
around active tracks so that initialisation concerns new objects only. The six camera layers of ViP3D are
replaced by two layers, one reading the \ac{bev} map at the reference point and one attending to the cameras
through spatially modulated attention, with the \ac{lidar} output kept for queries that no camera sees. The
model is trained in two stages, the \ac{lidar} model first and the camera layer on top of the frozen
\ac{lidar} feature extractor second. The realisation of this design is described in the following chapter
